\section{A quantitative Jacobi--Hardy column comparison}\label{sec:columns}
Continue with the angular notation \eqref{wang}. Put $\gamma=2/3$,
$N_j=j+3/2$, $c_0=11\pi/12$, and
\[
 v_j^\pm(\theta)=e^{\pm i(N_j\theta-c_0)}/\sqrt{2\pi},\quad
 v_j=v_j^-+v_j^+,\quad Ue_j=u_j,\quad V_\pm e_j=v_j^\pm,\quad V=V_-+V_+.
\]
These synthesis maps are bounded (the wave maps are restricted circle
Fourier synthesis). Set $X=M_{s_0}HU$, $X_0=M_{s_0}V_-$,
$A=X^*X$, $T_0=X_0^*X_0$.

\begin{theorem}\label{thm:columns}
One has $\sqrt j\|(A-T_0)e_j\|_2\to0$.
After restoring the exact basis phase $i^j$, the Toeplitz model has
real symmetric entries $b_{i-k}$ where
\begin{equation}\label{Toeplitzsymbol}
 b_r=\frac1{2\pi}\int_{-\pi/2}^{\pi/2}w(x+\pi/2)e^{irx}\,dx,
 \qquad |b_r|\le C(1+|r|)^{-5/3}.
\end{equation}
For the original $\psi$-coordinate matrix of $\mathcal B$, denoted
$B$, and this Toeplitz matrix $T_b$, one therefore has
\begin{equation}\label{originalcolumns}
 \|(B-T_b)e_j\|=o((j+1)^{-1/2}).
\end{equation}
\end{theorem}

\subsection{An endpoint Fourier lemma}
Suppose $k(\theta,\phi)$ is locally absolutely continuous in $\phi$,
with envelope $a_0(\theta)$ such that
\[
 |k|\le Ca_0(\theta)d(\phi)^{-1/2},\qquad
 |\partial_\phi k|\le Ca_0(\theta)d(\phi)^{-3/2}.
\]
Suppose that for each fixed interior $\theta$, the two scaled
expressions $d(\phi)^{1/2}k$ and $d(\phi)^{3/2}\partial_\phi k$
tend to zero at both endpoints. Then
\begin{equation}\label{Fourierlittleo}
 \sqrt N\int_0^\pi k(\theta,\phi)e^{\pm iN\phi}\,d\phi\longrightarrow0
\end{equation}
pointwise, bounded in absolute value by $Ca_0(\theta)$.
Split at $1/N$ and $\pi-1/N$. The endpoint pieces are
$O(a_0/\sqrt N)$; one integration by parts on the middle gives
the same order from its boundary and derivative terms. For little-o,
first choose fixed small endpoint neighborhoods where the scaled
bounds are arbitrarily small for this $\theta$, and integrate by
parts on the remaining compact interval, giving $O(1/N)$ there.
This proves \eqref{Fourierlittleo} for real $N\to\infty$.
Dominated convergence gives its $L^2(q(\theta)d\theta)$ version
whenever $\int q a_0^2<\infty$.

\subsection{Two explicit Hardy kernels}
Let $f(\theta)=\log\tan(\theta/2)$, so $f'=1/\sin\theta$.
The difference between $H$ and the compressed ordinary
negative-frequency Hardy operator $C_I$ has, apart from a fixed
factor $-i/(2\pi)$, the kernel
\[
 d_H(\theta,\phi)=
 \frac{\sqrt{f'(\theta)f'(\phi)}}{f(\theta)-f(\phi)}
                         -\frac1{\theta-\phi}.
\]
The canceled diagonal is filled continuously. The explicit bounds are
\begin{equation}\label{Hardyest}
 |d_H|\le\frac C{\sqrt{d(\theta)d(\phi)}},\qquad
 |\partial_\phi d_H|\le\frac C{\sqrt{d(\theta)}d(\phi)^{3/2}}.
\end{equation}
For completeness, in comparable left-endpoint scales set $\phi=\theta r$,
$1/2\le r\le2$. After extracting the scale, the canceled quotient
is smooth through $r=1$, and its required derivatives are uniformly
bounded; hence the bounds are $C/\theta$ and $C/\theta^2$.
In noncomparable scales use
$|f(\theta)-f(\phi)|\ge c|\log(\theta/\phi)|$, whose logarithm
is bounded away from zero. The ordinary Cauchy terms obey
\eqref{Hardyest} separately there. At opposite endpoints the denominator
is bounded below by
$c[1+|\log d(\theta)|+|\log d(\phi)|]$; mixed compact/endpoint
regions obey the same estimates. Direct differentiation uses
\[
 \partial_\phi\frac{\sqrt{f'(\theta)f'(\phi)}}{f(\theta)-f(\phi)}
 =\frac{\sqrt{f'(\theta)f'(\phi)}}{f(\theta)-f(\phi)}
       \left[-\frac{\cot\phi}{2}
                    +\frac{\csc\phi}{f(\theta)-f(\phi)}\right].
\]
For fixed interior $\theta$, the logarithmic denominator gives the
two endpoint little-o conditions. Since $\int w/d$ and
$\int w^2/d$ are finite, the Fourier lemma proves
\begin{equation}\label{Hardywaves}
 \sqrt j\|s_0(H-C_I)v_j^\pm\|_2\to0,
 \qquad \sqrt j\|w(H-C_I)v_j^\pm\|_2\to0.
\end{equation}

The commutator $[H,w]$ has, up to the same fixed factor, kernel
\[
 k_w(\theta,\phi)=(w(\phi)-w(\theta))
       \frac{\sqrt{f'(\theta)f'(\phi)}}{f(\theta)-f(\phi)}.
\]
It obeys the Fourier lemma with square-integrable envelope
\begin{equation}\label{commutenvelope}
 a_0(\theta)=C\left\{d(\theta)^{\gamma-1/2}
       +\frac1{\sqrt{d(\theta)}(1+|\log d(\theta)|)}\right\}.
\end{equation}
Here $|w^{(r)}|\le C_rd^{\gamma-r}$ for $r=0,1,2$.
At comparable endpoint scales the canceled quotient is
$O(\theta^{\gamma-1})$ and its derivative is
$O(\theta^{\gamma-2})$. For $\phi\ge2\theta$ near the same
endpoint, use
\[
 \sup_{2\theta\le\phi\le\eps}
 \frac{\phi^\gamma+\theta^\gamma}{\log(\phi/\theta)}
 \le C\left\{\theta^\gamma+(1+|\log\theta|)^{-1}\right\}.
\]
This follows by putting $v=\log(\phi/\theta)$; the term
$\theta^\gamma e^{\gamma v}/v$ has no interior maximum.
The case $\phi\le\theta/2$ follows directly from
$|w(\phi)-w(\theta)|\le C\theta^\gamma$. Opposite endpoints use
the logarithmic sum denominator from \eqref{Hardyest}; mixed
regions give the second summand in \eqref{commutenvelope}.
Differentiation adds $d(\phi)^{-1}$, using the displayed derivative
formula and the derivative bounds for $w$. For fixed $\theta$ the
scaled expressions still tend to zero, so
\begin{equation}\label{commuterate}
 \sqrt j\|[H,w]v_j^-\|_2\to0.
\end{equation}
The borderline envelope square $d^{-1}/(1+|\log d|)^2$ is integrable.

\subsection{The direct wave defect}
The normalized Jacobi expansion proved from the uniform Bessel
formula in \cite[Section 9]{r124} implies, for $R_j=u_j-v_j$,
\begin{equation}\label{Rbounds}
 |R_j(\theta)|\le C\min\{1,[(j+1)d(\theta)]^{-1}\},\quad
 \|R_j\|_2=O(j^{-1/2}),\quad
 \|s_0R_j\|_2=O(j^{-(1+\gamma)/2}).
\end{equation}
On every compact interior interval $\sqrt jR_j$ tends uniformly
to zero; its global $L^2$ norm is bounded. It is therefore weakly
null. The earlier proof shows $[s_0,H]$ compact, so
\[
 \sqrt j\|s_0HR_j\|_2\to0
\]
by $s_0HR_j=H(s_0R_j)+[s_0,H]R_j$.
This uses compactness on a bounded weakly null sequence with the
scale already included.
The ordinary compressed Hardy tail estimate is
\[
 |C_Iv_j-v_j^-|\le C\{L(N_j\theta)+L(N_j(\pi-\theta))\},\qquad
 L(r)=\begin{cases}1+|\log r|,&r\le1,\\1/r,&r\ge1.\end{cases}
\]
Integration against $w=O(d^\gamma)$ makes its weighted norm
$O(j^{-(1+\gamma)/2})$. Together with \eqref{Hardywaves} and
\eqref{Rbounds}, this proves
\begin{equation}\label{directdefect}
 \sqrt j\|(X-X_0)e_j\|_2\to0.
\end{equation}

\subsection{The adjoint defect must be controlled separately}
We prove
\begin{equation}\label{dualdefect}
 \sqrt j\|(U-V)^*(wv_j^-)\|_{\ell^2}\to0.
\end{equation}
First, uniformly in $k$,
\begin{equation}\label{variation}
 \Var(wR_k)\le C,\qquad (wR_k)(0)=(wR_k)(\pi)=0.
\end{equation}
Here no uncontrolled remainder is differentiated. The exact angular
Jacobi derivative identity for parameters $\alpha_J,\beta_J$ is
\[
 u_k'=q_J(\theta)u_k-
       \sqrt{k(k+\alpha_J+\beta_J+1)}\,
                      u_{k-1}^{(\alpha_J+1,\beta_J+1)},
\]
where $q_J=[(2\alpha_J+1)\cot(\theta/2)
-(2\beta_J+1)\tan(\theta/2)]/4$.
Apply the separate uniform first-order expansions to the two families.
The shifted family has the same $N_k$ and phase $\Phi_k-\pi/2$;
the leading derivative terms cancel, while
$\sqrt{k(k+3)}-N_k=O(1/N_k)$. Consequently
\[
 |R_k'|\le C[d^{-1}+(N_kd^2)^{-1}]\quad(N_kd\ge1),\qquad
 |R_k'|\le CN_k\quad(N_kd<1).
\]
The latter uses the exact identity and endpoint powers $11/6,7/6$,
both greater than one. Integrating $w'R_k+wR_k'$ with
\eqref{Rbounds} proves \eqref{variation}.

The first Jacobi correction is
\[
 A_J(\theta)=\frac{(1/4-\alpha_J^2)\cot(\theta/2)
                  -(1/4-\beta_J^2)\tan(\theta/2)}4,
 \quad (\alpha_J,\beta_J)=(4/3,2/3),\quad\Phi_k=N_k\theta-c_0.
\]
The normalized expansion, with no $1/k$ amplitude term, gives
\begin{equation}\label{L1first}
 \|w[kR_k-\sqrt{2/\pi}A_J\sin\Phi_k]\|_1\longrightarrow0.
\end{equation}
Indeed $wA_J\in L^1$. On $d<1/k$ both terms have integral
$O(k^{-\gamma})$; on the complement the uniform
$O((N_kd)^{-2})$ remainder, multiplied by $kw$, has the same bound.
Replacing $k/N_k$ by one adds $O(k^{-1})$.

Set $\beta_{kj}=\langle R_k,wv_j^-\rangle$. Integration by parts
in \eqref{variation} gives $|\beta_{kj}|\le C/(j+1)$ uniformly
in $k$. By \eqref{L1first}, $k\beta_{kj}$ is, uniformly in $j$,
a fixed linear combination of Fourier coefficients of $wA_J$ at
frequencies $k-j$ and $-(k+j+3)$, plus an error tending to zero
as $k\to\infty$. Fix $\eps>0$ and $M$. The contribution to
$j\sum_k|\beta_{kj}|^2$ from $k<\eps j$ is at most $C\eps$.
The window $|k-j|\le M$ contributes $O(M/j)$.
On $k\ge\eps j$ outside that window, Riemann--Lebesgue and the
uniform error in \eqref{L1first}, together with
$\sum_{k\ge\eps j}k^{-2}=O((\eps j)^{-1})$, give a limsup
bounded by $C\eps^{-1}\sup_{|r|>M}|\widehat{wA_J}(r)|^2$.
Take $j\to\infty$, then $M\to\infty$, then $\eps\downarrow0$.
This proves \eqref{dualdefect}.

Similarly $w$ vanishes at both endpoints and $w'\in L^1$.
Integration by parts and Riemann--Lebesgue give Fourier coefficients
$o(1/|r|)$ uniformly in the tail. The entries of $V_+^*(wv_j^-)$
have frequencies $k+j+3$, so summing their squares gives
\begin{equation}\label{plusdefect}
 \sqrt j\|V_+^*(wv_j^-)\|_{\ell^2}\to0.
\end{equation}
The exact decomposition
\[
 (H-I)wv_j^-=[H,w]v_j^-+w(H-C_I)v_j^-+w(C_I-I)v_j^-
\]
and \eqref{Hardywaves}--\eqref{commuterate} yield
\begin{equation}\label{weightedprojection}
 \sqrt j\|(H-I)(wv_j^-)\|_2\to0.
\end{equation}
The last ordinary Hardy term is $O(1/j)$ in $L^2$ because $2\gamma>1$.

Now put $D=X-X_0$. Exactly
\[
 A-T_0=X^*D+D^*X_0,
\]
\[
 D^*X_0e_j=U^*(H-I)(wv_j^-)
       +(U-V)^*(wv_j^-)+V_+^*(wv_j^-).
\]
Equations \eqref{directdefect}, \eqref{dualdefect},
\eqref{plusdefect}, and \eqref{weightedprojection} prove the first
claim of Theorem~\ref{thm:columns}.

\subsection{Returning to the original operator and restoring phase}
Let $S$ be unitary scaling to $z=y/(3\sqrt3)$, $\mathcal F$ the
unitary Fourier transform with kernel $e^{-iz\xi}$, and $J$ the
Fourier-to-angular coordinate unitary. With the scaled interpretation
of $v$, the exact transform $L=J\mathcal FS M_v$ satisfies
\[
 Lf_j=i^ju_j,\qquad
 J\mathcal FS C_0S^*\mathcal F^*J^*=HM_wH.
\]
The latter identity has no gamma phase. Therefore
\[
 L\widetilde{\mathcal B}f_j-(HM_wH)Lf_j
   =L\mathcal Ef_j+J\mathcal FS(M_vC_0-C_0M_v)f_j.
\]
The second term is the unitary image of
$(M_vC_0M_v^*-C_0)M_vf_j$. The two precise smoothing variants
in Theorem~\ref{thm:smoothing} prove an $o(j^{-1/2})$ error.
Let $Q=\diag(i^j)$. The original coordinate matrix is
$Q^*U^*L\widetilde{\mathcal B}L^*UQ$, and the Hardy model is
$Q^*AQ$. Column norms are unchanged by the diagonal phases.
The raw wave Gram entry is
$(2\pi)^{-1}\int_0^\pi w(\theta)e^{i(i-k)\theta}\,d\theta$.
Conjugating by $Q$ multiplies it by $i^{k-i}$; setting
$x=\theta-\pi/2$ proves exactly \eqref{Toeplitzsymbol} and
\eqref{originalcolumns}.

The corresponding circle symbol is $w(x+\pi/2)$ on
$|x|\le\pi/2$ and zero elsewhere. It is endpoint $C^{2/3}$,
with first and second derivative bounds $d^{-1/3},d^{-4/3}$.
Integrate once by parts, split at distance $1/|r|$ from each
support endpoint, and integrate the interior part once more by
parts. This proves $b_r=O(|r|^{-5/3})$.
The symmetry $w(\pi-\theta)=w(\theta)$ makes $b_r$ real and even.
The absolutely convergent Fourier series vanishes at $x=\pi$, so
\begin{equation}\label{alternatingsum}
 \sum_{r\ge1}(-1)^rb_r=-b_0/2=\sqrt3/4-1/2,
 \qquad b_0=1-\sqrt3/2.
\end{equation}
The last value is the exact scalar mean already evaluated in
\cite[Section 9]{r124}. All parts of Theorem~\ref{thm:columns}
are now proved.
